\section{Limitations and Scope}\label{sec:open}

The paper's central object is the temporal-validity horizon for finite-memory
tracking under drift. The main closed results are the general horizon law, the
uniform-window staleness bound, the one-dimensional root-$n$ proof model, the
structural Gaussian lower bound at the exponent level, and the Gaussian
location minimax benchmark. Other parts of the manuscript are intentionally
stated with narrower status.

\paragraph{Finite-sample scope.}
The theory is modular in the finite-sample exponent $a$, but the main rigorous
instantiation is the root-$n$ regime. The tractable proof model establishes that
regime in one dimension under bounded support and fixed span. Outside that
setting, the same balancing logic continues to hold, but the finite-sample term
must come from a different theorem. In particular, the paper does not claim a
general raw-$W_2$ triangular-array theorem in arbitrary dimension.

\paragraph{Operational scope.}
For fixed-$\varepsilon$ Sinkhorn, the manuscript proves the structural
ingredients of the embedded fixed-span model and states the full triangular-array
inheritance law as a conjecture. Exact support-complexity inheritance and the
dual threshold $\alpha>k/2$ are closed; the remaining step is the full theorem
connecting the triangular-array finite-sample term to the i.i.d. benchmark.

\paragraph{Family-level scope.}
The manuscript does not claim a blanket extension to every deterministic
$W_2$-H\"older path class. Continuous regular families and operational
geometries can support meaningful finite-sample exponents, but support-changing
and singular families behave differently. The natural extension target is
therefore a theorem for regular dominated families or operational geometries
with comparable local testing structure.

\paragraph{Lower-bound scope.}
The structural lower bound is subclass based. The Gaussian construction is
enough to show that, in the root-$n$ regime, the horizon is structural rather
than an artifact of a particular estimator. What remains open is class-tightness
and constant-sharpness for broader deterministic H\"older path classes and for
finite-sample regimes beyond $a=1/2$. The refined Gaussian geometry sharpens the
constants available within that lower-bound program, but it does not close the
full distributional lower theory.

\paragraph{Weights and memory shape.}
The main text emphasizes uniform windows. The paper does not investigate
whether broader non-uniform weights can improve constants or alter exponents
across richer path classes.

\paragraph{Online roughness estimation.}
The law is stated in terms of the roughness parameters $H$ and $\zeta$, but the
paper does not propose a fully online procedure for estimating them or for
quantifying sensitivity to misspecification. Turning the horizon law into a
robust adaptive policy remains future work.

\paragraph{Empirical scope.}
The experiments are synthetic and structurally motivated. Their role is to make
the theory visible and to provide evidence for the operational conjecture, not to demonstrate
competitive performance on a benchmark suite. Real-data validation and broader
comparisons with adaptive horizon-selection rules remain outside the present
scope.

\paragraph{Detectability versus validity.}
Validity loss before detector alarm is an empirical consequence of the theory,
not the primary theorem object. The paper shows that temporal validity can fail
before detector alarms appear, but it does not prove a universal detection-delay
theorem across detector classes.

\paragraph{Open problems.}
The next technical questions are clear. The first is to sharpen the
finite-sample constant in the one-dimensional proof model. The second is to
understand when practical geometries, such as fixed-$\varepsilon$ Sinkhorn,
admit full triangular-array inheritance theorems in genuinely mid- or
high-dimensional settings. The third is to determine which regular dominated
distribution families support which finite-sample exponents and to extend the
lower theory accordingly. Beyond that come class-tight lower bounds,
non-uniform weighting, real-data validation, and stronger stochastic-process
formulations.
